% =====================================================================
% Strand Two: the geometry of the internal Self/World representations.
% Intended to slot into the "Analysis of Trained Networks" section,
% expanding the "Belief geometry" and "Self-recognition" paragraphs.
% Reuses only citation keys already present in references.bib.
% =====================================================================

\subsection*{The geometry of the Self/World factorisation}

The analyses above ask \emph{whether} the SWITCH factorisation is recovered; here we
make precise \emph{what geometric object} we expect to find, so that the probes and
interventions have a concrete target. In the passive predictor of
\citep{shai2024belief} the object is well understood: a belief is a point in the
probability simplex $\Delta_W$ over the World hidden states, the Bayes-optimal belief
given a history is a single point in $\Delta_W$, and the set of all reachable beliefs
is the (typically fractal) mixed-state presentation (MSP), the attractor of the
iterated function system whose maps are the observation-conditioned Bayes updates. Two
computations sit on top of this geometry: a \emph{linear} emission read-out giving the
next-token distribution, and a \emph{nonlinear} filtering update that moves the belief
point as each observation arrives, both of which \citep{shai2024belief} recover
linearly from the residual stream.

\paragraph{From one simplex to a product.} The SWITCH hidden state is the pair $(s,w)$,
so the belief lives in the joint simplex over $S\times W$. The central hypothesis,
stated geometrically, is that the reachable-belief manifold collapses onto the
\emph{product} submanifold $\Delta_S\times\Delta_W$ --- that the joint belief factors as
$b(s,w)\approx b_S(s)\,b_W(w)$ with the two marginals independently decodable --- rather
than filling an entangled subset of the joint simplex. This is the precise sense in
which Self and World are ``separable'', and it ties our result to the Simplex line of
work on the representation of independent generative processes. Token parity now selects
between two semantically distinct updates rather than one: an \emph{observation} token
triggers a filtering update that contracts $b_W$ toward consistency with what was seen,
while an \emph{action} token triggers a \emph{control} pushforward
$b_W \mapsto P(\cdot\mid\cdot,a)^{\top} b_W$ that advances the World belief through the
action-conditioned dynamics \citep{bengio1995iohmm}. Recovering two distinct update
maps, selected by parity and acting on a product geometry, is the first concrete target
of this strand.

\paragraph{The Self simplex as a persona-selection model.} In Regimes 2--4 the model is
always the actor, so $b_S$ is a point mass: there is no uncertainty about who is acting
and the Self marginal is degenerate, leaving only the thin token-type (action vs.\
observation) distinction. The Self simplex becomes a non-trivial object precisely in the
mixed-actor Regime 5, where each trajectory is generated by one of a family of policies
$\{\pi^{(1)},\dots,\pi^{(M)}\}$, one of which ($\pi^{(0)}$) is the model itself. Here
$\Delta_S$ is the online posterior over \emph{actor identity}: its vertices are the
members of the actor family, and it is updated by the likelihood ratio of each observed
action under the candidate policies --- exactly a filtering process on a ``which-actor''
latent. We therefore expect a second, lower-dimensional mixed-state geometry --- a
persona-selection MSP --- to appear alongside the World MSP, linearly decodable from the
residual stream in Regime 5 but absent in Regimes 2--4.

\paragraph{Nested belief and theory of mind.} The persona simplex does more than label
trajectories: it gates the World update. An action is evidence about the World state
only insofar as the acting policy is World-state-dependent, and how strongly depends on
which actor produced it. The action-token update on $b_W$ therefore carries an
actor-dependent Bayesian factor \emph{before} the pushforward,
\[
  b_W(w)\;\propto\; b_W(w)\,\pi^{(c)}(a\mid w)\quad\text{(actor-gated)},
  \qquad\text{then}\qquad
  b_W \;\mapsto\; P(\cdot\mid\cdot,a)^{\top} b_W,
\]
so that a random-actor action is treated as uninformative (pushforward only) while an
optimal-actor action drives a strong belief update. When the actors are themselves
\emph{stateful} --- e.g.\ Bayes-adaptive agents carrying their own World belief ---
predicting their actions requires the model to simulate \emph{that actor's} belief, a
separate point in a belief simplex that need not coincide with the model's own. The full
latent geometry is then a \emph{bundle}: over each non-self node $c$ of $\Delta_S$ sits a
fiber simplex $\Delta_W^{(c)}$ encoding the model's representation of actor $c$'s
beliefs, distinct both from the ground-truth World state and from the model's own
belief. Representing another agent's belief as separable from one's own is first-order
theory of mind; the SWITCH with stateful actors forces it as a computational requirement
of observation prediction, and it is jointly testable with the World geometry on the
same analytic ground truth.

\paragraph{The self as a degenerate fiber.} The self node is distinguished not by its
location in $\Delta_S$ but by the status of its fiber: over the self node the belief
fiber is \emph{identified} with the model's own World belief $b_W$ rather than separately
inferred, because the self acts on its own belief by construction. This identification is
the geometric content of the reflexive self-recognition reported behaviourally in
post-trained models \citep{betley2025tellme,panickssery2024self} and in our forthcoming
work with Lindsey. It is also what makes the Self component causally load-bearing: a
perturbation written into the self fiber should move both the model's observation
predictions and its own action distribution, whereas a perturbation to a non-self fiber
$\Delta_W^{(c)}$ should move only the predicted behaviour of actor $c$. This double
dissociation, layered on top of the Self-steering interventions already described, is the
sharpest causal signature of a genuine self representation.

\paragraph{Circuit-level predictions and the off-policy test.} Two implementations are
distinguishable and worth separating experimentally. Under a \emph{shared, identity-gated}
update circuit, a single mechanism performs the actor-conditioned belief update with the
$\Delta_S$ vector entering as a gating input; steering $\Delta_S$ should then continuously
deform the observed update rule, from ``ignore the action as evidence'' to ``treat it as
strong evidence''. Under a \emph{privileged self pathway}, the self is handled by the
first-person generation circuit --- read off, not inferred --- rather than by the shared
theory-of-mind module, and only the non-self nodes route through the latter.
Distinguishing these is the mechanistic content of the off-policy generalisation test
already proposed: a shared identity-gated bundle should interpolate to a held-out actor
of a known type (Hypothesis A, generalisation), whereas per-actor shortcuts through the
joint (action, observation) statistics should not (Hypothesis B). Two design notes follow
directly: the nested fiber structure arises only if the actor family contains stateful
policies (memoryless actors require actor identity but no fiber), and the actors should
remain first-order World modellers in the toy --- recursive actor-models-actor structure
is the layered self-simulation we defer to the language-model extension.
